\subsection{与埃尔米特型相伴的双线性型。}

这里仍假设环 A 是交换环 K 的二次扩张 \(\mathrm{A} = \mathrm{K}(i)\)（其中 \(i^{2} = -1\)），且 J 是 A 的自同构 \(\lambda + i\mu \to \lambda - i\mu\)（ \(\lambda \in K, \mu \in K\)）。设 E 与 \(E'\) 是两个 A-模，\(E_{0}\) 与 \(E_{0}'\) 是 E 与 \(E'\) 的底 K-模，j 与 \(j'\) 是自同构 \(x \to ix\) 与 \(x' \to ix'\)（E 与 \(E'\) 的）（参见第 2 小节）。称 \(E_{0} \times E_{0}'\) 上的 K-双线性型 f 在 j 与 \(j'\) 下不变，如果有

\[
f (j (x), j ^ {\prime} (x ^ {\prime})) = f (x, x ^ {\prime})\tag{5}
\]

对 \(x \in \mathrm{E}_0\) 与 \(x' \in \mathrm{E}_0'\)。把 \(x\) 换成 \(j(x)\)，可见该条件等价于

\[
f (x, j ^ {\prime} (x ^ {\prime})) = - f (j (x), x ^ {\prime})\tag{6}
\]

对一切 \(x \in \mathbf{E}_0\) 与 \(x' \in \mathbf{E}_0'\)。

命题 1. --- 设 \(\Phi_{1}\)（相应地 \(\Phi_{2}\)）是 \(\mathbf{E}_{0} \times \mathbf{E}_{0}^{\prime}\) 上在 \(j\) 与 \(j^{\prime}\) 下不变的 K-双线性型。使 \(\Phi_{1}\)（相应地 \(\Phi_{2}\)）对应于映射 \(\Phi\) （\(\mathbf{E} \times \mathbf{E}^{\prime}\) 到 A 中的、由下式定义的）的那个映射

\begin{enumerate}
\def\labelenumi{(\arabic{enumi})}
\setcounter{enumi}{6}
\tightlist
\item
\end{enumerate}

\[
\Phi (x, x ^ {\prime}) = \Phi_ {1} (x, x ^ {\prime}) + i \Phi_ {1} (x, j ^ {\prime} (x ^ {\prime}))\tag{8}
\]

\[
(\text { resp. } \Phi (x, x ^ {\prime}) = - \Phi_ {2} (x, j ^ {\prime} (x ^ {\prime})) + i \Phi_ {2} (x, x ^ {\prime}))
\]

\((x \in \mathrm{E}, x' \in \mathrm{E})\)，是 \(\mathrm{E}_0 \times \mathrm{E}_0'\) 上在 \(j\) 与 \(j'\) 下不变的 K-双线性型所成的 K-向量空间到 \(\mathrm{E} \times \mathrm{E}'\) 上半双线性型所成的 K-向量空间上的一个同构。再设 \(\mathrm{E} = \mathrm{E}'\)；要使 \(\Phi\) 是埃尔米特型，必须且只需 \(\Phi_1\) 是对称的（相应地 \(\Phi_2\) 是反对称的）（参见习题 4）。

事实上，E × E' 到 A 中的每个映射 \(\Phi\) 都可以以唯一的方式写成 \(\Phi = \Phi_1 + i\Phi_2\) 的形式，其中 \(\Phi_1\) 与 \(\Phi_2\) 是 E × E' 到 K 中的映射。由第 2 小节公式 (3)（相应地 (4)），要使偏映射 \(x \to \Phi(x, x')\) 是 A-线性的，必须且只需 \(\Phi_1\)（相应地 \(\Phi_2\)）关于 \(x\) 是 K-线性的，且有

\begin{enumerate}
\def\labelenumi{(\arabic{enumi})}
\setcounter{enumi}{8}
\tightlist
\item
\end{enumerate}

\[
\Phi (x, x ^ {\prime}) = \Phi_ {1} (x, x ^ {\prime}) - i \Phi_ {1} (j (x), x ^ {\prime})\tag{10}
\]

\[
(\text { resp. } \Phi (x, x ^ {\prime}) = \Phi_ {2} (j (x), x ^ {\prime}) + i \Phi_ {2} (x, x ^ {\prime})).
\]

类似地，要使 \(\overline{\Phi(x,x^{\prime})}\) 关于 \(x^{\prime}\) 是 A-线性的，必须且只需 \(\Phi_{1}\)（相应地 \(\Phi_{2}\)）关于 \(x^{\prime}\) 是 K-线性的，且有

\begin{enumerate}
\def\labelenumi{(\arabic{enumi})}
\setcounter{enumi}{10}
\tightlist
\item
\end{enumerate}

\[
\Phi (x, x ^ {\prime}) = \Phi_ {1} (x, x ^ {\prime}) + i \Phi_ {1} (x, j ^ {\prime} (x ^ {\prime}))\tag{12}
\]

\[
(\text { resp. } \Phi (x, x ^ {\prime}) = - \Phi_ {2} (x, j ^ {\prime} (x ^ {\prime})) + i \Phi_ {2} (x, x ^ {\prime})).
\]

由此立即得出，要使 \(\Phi\) 是半双线性型，必须且只需它可以写成 (9) 与 (11)（相应地 (10) 与 (12)）两种形式之一，且 \(\Phi_{1}\)（相应地 \(\Phi_{2}\)）是在 \(j\) 与 \(j'\) 下不变的 K-双线性型。

由此可知，要使半双线性型 \(\Phi = \Phi_{1} + i\Phi_{2}\) 为零，必须且只需 \(\Phi_{1}\)（相应地 \(\Phi_{2}\)）为零。而若 \(E = E'\)，则 \(\Phi(y, x) = \Phi_{1}(y, x) + i\Phi_{2}(y, x)\) 且 \(\overline{\Phi(x, y)} = \Phi_{1}(x, y) - i\Phi_{2}(x, y)\)；要使这两个表达式相等，换言之要使 \(\Phi\) 是埃尔米特型，必须且只需 \(\Phi_{1}\) 对称（相应地 \(\Phi_{2}\) 反对称）。

注. --- 1) 公式 (7) 与 (8) 表明，若 \(x \in E\)，要使 \(\Phi(x, x') = 0\) 对一切 \(x' \in E'\) 成立，必须且只需 \(\Phi_1(x, x') = 0\)（相应地 \(\Phi_2(x, x') = 0\)）对一切 \(x' \in E'\) 成立。

\begin{enumerate}
\def\labelenumi{\arabic{enumi})}
\setcounter{enumi}{1}
\tightlist
\item
  E 的自同态 \(u\) 关于 \(\Phi\) 的伴随（\(\S 1\)，第 8 小节）与 \(u\)（视为 \(\mathrm{E}_{0}\) 的自同态）关于 \(\Phi_{1}\)（相应地 \(\Phi_{2}\)）的伴随相同。
\end{enumerate}

